KIRIN est un projet pédagogique de développement d'un solveur
bidimensionnel d'écoulements incompressibles, écrit en Rust.
Son objectif principal est de présenter l'ensemble du cheminement
qui conduit des équations physiques à leur implémentation :
choix des hypothèses, formulation du modèle, discrétisation,
construction des algorithmes et traduction en code.

Deux approches numériques seront explorées :
\begin{itemize}
    \item une approche conservative, privilégiant la fidélité
    des résultats et le respect des bilans de conservation ;

    \item une approche inspirée de \textit{Stable Fluids}
    de Jos Stam (1999)~\cite{stam1999stable}, privilégiant la robustesse et l'interactivité.
\end{itemize}

Ce document accompagne cette démarche en explicitant les choix
physiques et numériques, leur mise en œuvre et la vérification
des résultats obtenus.
